% !TeX root = ../Trabajo_RPM.tex
\subsection{Interpretación económica y dimensional}
La función objetivo asigna el mismo peso a un minuto en cualquier máquina y mide la carga agregada de trabajo. No incorpora diferencias de consumo eléctrico, personal o tarifas entre equipos. En el caso de añadir costos por minuto el objetivo adquiriría un carácter monetario y los coeficientes se verían alterados; los óptimos que se dan no son directamente interpretables bajo ese nuevo objetivo.

Como nota, dividir lotes permite economías de procesamiento, ya que se pueden utilizar máquinas de mayor velocidad, pero se suele realizar a costa de un mayor número de activaciones. Esto se refleja en el modelo por $C_{ij}x_{ijt}+S_jy_{ijt}$. No se trata de tener siete productos cuando tenemos 207 activaciones en tres días con 200 referencias, una sola referencia puede provocar que tengamos más de una activación extra. El número de referencias activadas varias veces y el número de referencias activadas en varias máquinas se da en la Sección~\ref{sec:resultados}.
Por otro lado, la restricción de capacidad es local. Es necesario que se cumpla $Z_H\leq\sum_{j,t}K_{jt}$ pero no es suficiente. Podríamos sobrepasar la capacidad de una de las máquinas aunque no nos pasemos en el resto de las máquinas. Por esta razón se considera cada combinación máquina y día en lugar del promedio de utilización. En cierto sentido, esta propiedad es la que complica el problema y por este motivo se ha de verificar la distribución temporal del maestro agregado.

\subsection{Por qué importa la cota de enlace}
Al relajar el problema, la restricción $x_{ijt}\leq M y_{ijt}$ implica que la variable $y_{ijt}$ debe cumplir $y_{ijt}\geq x_{ijt}/M$. Si $M$ es grande se puede disminuir el valor de $y_{ijt}$ hasta ese valor fraccionario y únicamente se pagaría $S_jx_{ijt}/M$ por la preparación si se puede hacer cumpliendo con el resto de las restricciones. Sin embargo, en el problema entero, si hay producción, se debe tener el valor completo de $y_{ijt}$. Esta diferencia puede implicar que se tenga una peor cota inferior, pero no cambia el óptimo entero.

La cota mostrada en la ecuación~\eqref{eq:U} se obtiene utilizando dos restricciones que deben ser cumplidas por cualquier solución factible: no puede producirse más de lo que se demanda, y no puede utilizarse más de la capacidad disponible una vez hecha la preparación. Esto mantiene el conjunto de soluciones enteras del modelo original, y no se han hecho cambios como el modificar los tiempos ni el redondear la demanda para hacer más fácil la prueba. La variable entera no puede tomar valores mayores que la capacidad (que en el problema relajado es fraccionaria) de unidades, por eso se impone el piso.

El corte de cobertura viene a completar el enlace. Si bien la igualdad de la demanda y los enlaces obligan a tener al menos una activación en una solución entera donde $D_i>0$, la relajación admite activaciones fraccionarias repartidas entre varios pares. La condición $\sum_{j,t}y_{ijt}\geq1$ evita que se pague menos de una activación total por referencia. La mejora de la cota depende de la interacción entre los refuerzos; con los datos del experimento no se puede cuantificar su efecto por separado.

\subsection{Orden de las decisiones de horizonte y carga}
Reducir el tiempo agregado y localizar el horizonte más pequeño son dos objetivos distintos. El método para avanzar es una regla de tipo lexicográfico: se fija $H^*$, en la medida de lo posible, y luego se minimiza $Z_{H^*}$. No es cuestión de sumar un día y un minuto con pesos arbitrarios. Un horizonte mayor puede reducir la carga agregada, aunque deje de ser el menor plazo posible en días.

Dadas unas reglas de alistamiento y una capacidad diaria que se va repitiendo, todo plan viable de $H$ días puede extenderse a $H+1$ con el último exento de producción. De ello se deduce que $\mathcal F_H\subseteq\mathcal F_{H+1}$ si se asocia cada plan con su nula extensión de producción en el día de más, y que $Z_{H+1}^*\leq Z_H^*$. Esta propiedad parte del supuesto de que al ensanchar la ventana no se introducen nuevas demandas o restricciones, aunque no por ello se garantiza un descenso estricto del objetivo.

Si se quiere analizar una ventana más amplia que la ventana mínima demostrada, entonces, debe resolverse el modelo con el tamaño deseado de la ventana, pero manteniendo la misma demanda fija. En este caso, el certificado de carga es para la ventana analizada, y la prueba para el mínimo de días es la prueba de factibilidad. El horizonte automático es el que se utiliza en la versión de la aplicación descrita en este trabajo. La propiedad enunciada anteriormente es de la formulación y no es que la aplicación necesariamente tenga un selector del horizonte manualmente.